\section*{الفصل \ref{ch:b2:biogeochemical-cycles} --- الدورات البيوجيوكيميائية}

\begin{solution}{exo:b2:biogeochemical-cycles:1}
\omterm{def:b2:biogeochemical-cycles:reservoir}{المستودع} مخزن من العنصر (أي كتلة)؛ \omterm{def:b2:biogeochemical-cycles:reservoir}{والتدفق} معدل
انتقال بين \omterm{def:b2:biogeochemical-cycles:reservoir}{المستودعات} (كتلة في السنة)؛ \omterm{def:b2:biogeochemical-cycles:reservoir}{وزمن المكث}
كتلة \omterm{def:b2:biogeochemical-cycles:reservoir}{المستودع} مقسومة على خرجه الكلي، أي الزمن الوسطي الذي
تقضيه فيه ذرة. فنباتات اليابسة: $450/120 \approx 4$ سنوات (وهو تقريبًا
عمر كربون ورقة؛ أما \omterm{def:b2:plant-meristems:wood}{الخشب} فيبقى عقودًا ويخفي المتوسط
التشتت). والمحيط العميق: $37000/100 \approx 400$ سنة --- أي الزمن
قبل أن يرى جزيء هابط السطحَ من جديد.
\end{solution}

\begin{solution}{exo:b2:biogeochemical-cycles:2}
$2.5\times 2.13 = \qty{5.3}{GtC}$ في السنة. و$10/2.13 = \qty{4.7}{ppm}$
--- لو بقي كله في الهواء؛ وعند كسر محمول جوًا قدره \qty{45}{\%}،
\qty{2.1}{ppm}.
\end{solution}

\begin{solution}{exo:b2:biogeochemical-cycles:3}
% Long unbreakable runs (a chain of numerals, a Latin binomial) are one
% directional box under bidi=basic and cannot be broken; \sloppy must be in
% force when the paragraph is BROKEN, hence the \par before \endgroup.
\begingroup\sloppy
الفتح: التثبيت الحيوي بالنتروجيناز (الزراقم،
والريزوبيا، و\emph{Azotobacter}، و\emph{Frankia}) والبرق. والإعادة:
\omterm{def:b2:microbial-metabolism:anaerobic}{نزع النترتة} (لاهوائيات اختيارية، مثل
\emph{Pseudomonas} و\emph{Paracoccus}، في التربة
والرواسب اللاأكسجينية) والأناموكس
(البلانكتوميسيتات). وبينهما: أمنجة النتروجين
العضوي إلى أمونيوم بالمحلِّلات (الفطريات
والبكتيريا)؛ ونترتة
الأمونيوم إلى نتريت (\emph{Nitrosomonas}، والعتائق
المؤكسدة للنشادر) ونترتة النتريت إلى نترات (\emph{Nitrobacter})؛
واستيعاب الأمونيوم أو النترات بالنبات والمجهريات.
\par\endgroup
\end{solution}

\begin{solution}{exo:b2:biogeochemical-cycles:4}
على المدى الجيولوجي يمكن صنع النتروجين دائمًا من
$\mathrm{N_2}$ الهواء الذي لا ينفد بالمثبِّتات، وهي تُحبَّذ
كلما شحّ النتروجين؛ أما الفوسفور فلا \omterm{def:b2:biogeochemical-cycles:reservoir}{مستودع} كهذا له
ويقرر عرضَه التجويةُ، فيكون السقف بعيد المدى على
الإنتاجية هو الفوسفور. وفي موسم واحد يكون التثبيت بطيئًا
ومكلفًا، ومعظم النبات لا يقدر عليه، فيكون النتروجين المتاح
هو ما يحد النمو هنا والآن --- ولذلك كان السماد
نتروجينيًا في معظمه.
\end{solution}

\begin{solution}{exo:b2:biogeochemical-cycles:5}
$\tau = 1/0.2 = 5$ سنوات؛ و$M^{*} = F\tau = \qty{10}{t}$، أي
$\qty{1}{g/m^3} = \qty{1000}{mg/m^3}$، وهو إتخام فادح. والزمن لبلوغ
\qty{90}{\%}: $\tau\ln 10 = 11.5$ سنة. وتنصيف الدخل ينصّف
\omterm{def:b2:biogeochemical-cycles:reservoir}{الحالة المستقرة}، إلى \qty{500}{mg/m^3}، وتُبلغ على المقياس الزمني
نفسه.
\end{solution}

\begin{solution}{exo:b2:biogeochemical-cycles:6}
الستينيات: \qty{0.9}{ppm/yr}، أي \qty{1.9}{GtC/yr}. والعقد الثاني من هذا القرن: \qty{2.3}{ppm/yr}،
أي \qty{5.0}{GtC/yr}. فقد ارتفع معدل النمو مرتين ونصفًا،
على وتيرة الانطلاقات.
\end{solution}

\begin{solution}{exo:b2:biogeochemical-cycles:7}
\qty{6}{ppm} هي \qty{13}{GtC}، في مقابل إنتاج أولي صافٍ قدره
\qty{35}{GtC}: فسن المنشار ثلثه. وهو أصغر لأن
التنفس والتحلل يمضيان طوال الصيف (فسن المنشار
يسجل الالتقاط الصافي لا الإجمالي)، ولأن فصول نصف الكرة
الجنوبي معاكسة فتلاشي بعضه، ولأن
المحيط يخمد التأرجح.
\end{solution}

\begin{solution}{exo:b2:biogeochemical-cycles:8}
$\qty{200}{kg}/(\qty{28}{g/mol}) = 7100$~مول من $\mathrm{N_2}$؛ و$16\times
7100 = \num{114000}$~مول من ATP؛ و$/30 = 3800$~مول غلوكوز، أي \qty{690}{kg}.
والمحصول الذي يبني \qty{10}{t} من المادة الجافة يثبّت نحو
\qtyrange{15}{20}{t} من السكر إذا حُسب التنفس، فيكلفه
التثبيت نحو \qty{4}{\%} من ناتج تركيبه الضوئي --- وهو
ثمن الاستقلال عن نتروجين التربة.
\end{solution}

\begin{solution}{exo:b2:biogeochemical-cycles:9}
النسبة المتاحة $\mathrm{N:P} = 32/2.2 = 14.5$، دون نسبة ردفيلد
16: فتنفد النترات أولًا، ويبقى \qty{0.2}{\micro mol/kg} من
الفوسفات. والكربون المثبَّت: $32\times 106/16 =
\qty{212}{\micro mol/kg}$، أي نحو \qty{2.5}{mg} من الكربون في
الكيلوغرام من الماء.
\end{solution}

\begin{solution}{exo:b2:biogeochemical-cycles:10}
$E^{*} = Q\tau = \qty{500}{GtC}$، أي \qty{235}{ppm} فوق
285 ما قبل الصناعة: أي نحو \qty{520}{ppm}. وهو أسوأ في الواقع لأن
المصارف ليست حوضًا سريعًا واحدًا: فالمحيط السطحي يتشبع
(إذ تهبط سعته لالتقاط $\mathrm{CO_2}$ كلما استُهلك كربوناته)،
ويضعف مصرف اليابسة مع الاحترار، وللإزالة الحقيقية مركّبة
بطيئة --- أي امتزاج المحيط العميق وتجوية الصخر --- ذات
ثوابت زمنية من ألف إلى مئة ألف سنة، فلا يسترخي
كسر من الفائض عند أي $\tau$ بالعقود.
\end{solution}

\begin{solution}{exo:b2:biogeochemical-cycles:11}
ما قيس هو تبادل $\mathrm{CO_2}$ الجوي مع
المحيط والمحيط الحيوي، لا التفكك: فقد مُدّد $\mathrm{{}^{14}C}$ القنابل
في حوضي المحيط السطحي والنبات والتربة الأكبر بكثير،
وأعادا كربونًا عاديًا. وعمر النصف البالغ 11 سنة
للفائض (أي زمن تناقص أسّي قدره نحو 16 سنة) هو المقياس
الزمني لذلك التبادل --- أي دوران كربون الغلاف الجوي
في مقابل المحيط السطحي واليابسة، وهو أطول من
\omterm{def:b2:biogeochemical-cycles:reservoir}{زمن المكث} الإجمالي البالغ أربع سنوات لأن جلّ ما يغادر يعود
في الحال.
\end{solution}

\begin{solution}{exo:b2:biogeochemical-cycles:12}
النتروجين: التثبيت الطبيعي نحو 200~Tg في السنة؛ والتثبيت البشري
(هابر--بوش 120، والمحاصيل 60، والاحتراق 30) نحو 210: \omterm{def:b2:biogeochemical-cycles:reservoir}{فالتدفق}
مضاعف، \omterm{def:b2:biogeochemical-cycles:reservoir}{والمستودعات} المتأثرة --- التربة والأنهار والبحار
الساحلية و$\mathrm{N_2O}$ الهواء --- صغيرة سريعة، فيظهر التغير
في غضون عقود، في حين لا يُمَسّ \omterm{def:b2:biogeochemical-cycles:reservoir}{مستودع} $\mathrm{N_2}$.
أما الكربون: فالانطلاقات البشرية البالغة 11~GtC في السنة ليست إلا \qty{5}{\%} من
التبادل الطبيعي الإجمالي البالغ نحو 200، \omterm{def:b2:biogeochemical-cycles:reservoir}{فالتدفق} موكوز؛
لكن التدفقات الطبيعية تتلاشى \omterm{def:b2:biogeochemical-cycles:reservoir}{والتدفق} البشري لا يتلاشى، وقد
رفع \omterm{def:b2:biogeochemical-cycles:reservoir}{المستودع} الجوي بمقدار \qty{50}{\%} --- فالوكزة الصغيرة
لها أثر تراكمي كبير في \omterm{def:b2:biogeochemical-cycles:reservoir}{مستودع} مصرفه النهائي بطيء.
وأي الدورتين «أشد اضطرابًا» يتوقف على ما إذا
عُدّت التدفقات أم عُدّ التراكم.
\end{solution}

\begin{solution}{pb:b2:biogeochemical-cycles:1}
\textbf{1.} $285\times 2.13 = \qty{607}{GtC}$؛ و$412\times 2.13 =
\qty{878}{GtC}$؛ والزيادة \qty{271}{GtC}.
\textbf{2.} $271/700 = 0.39$ (والكسر على مدى العقود الأخيرة،
بحساب انطلاقات استعمال الأرض، أقرب إلى $0.45$).
\textbf{3.} المحيط، بنحو ربع الانطلاقات، واليابسة
(بنمو الغابات من جديد وبتسميد $\mathrm{CO_2}$)، بنحو الثلث؛ وهما
بقية \qty{430}{GtC} التي ليست في الهواء.
\textbf{4.} $175/38000 = \qty{0.5}{\%}$ --- وهو مهمل للمحيط
كله؛ لكن الالتقاط دخل حتى الآن الطبقة السطحية أساسًا
(\qty{900}{GtC})، وهو فيها ارتفاع بعُشر أو أكثر، وكيمياء
الكربونات تحوّل ارتفاعًا صغيرًا في $\mathrm{CO_2}$ المنحل إلى
ارتفاع أكبر في تركيز شوارد الهدروجين: فقد هبطت حموضة السطح
بمقدار $0.1$، أي ارتفاع في الحموضة قدره \qty{30}{\%}.
\textbf{5.} عند $Q = Q_0 e^{t/T}$، جرّب $E = Ae^{t/T}$: فيكون $A/T = Q_0 -
A/\tau$، ومن ثم $A = Q_0\tau T/(\tau + T)$ ويكون \omterm{prop:b2:biogeochemical-cycles:carbon}{الكسر المحمول جوًا}
$\dot E/Q = \tau/(\tau + T)$ ثابتًا. وانطلاقات تنمو بمقدار
\qty{2}{\%} في السنة ($T = 50$) عند $\tau = 40$ تعطي $0.44$: فالكسر
ثابت لأن المصرف والمنبع ينموان معًا.
\textbf{6.} عند $Q$ ثابت يؤول $E \to Q\tau$ ويؤول $\dot E \to 0$: فيهبط
\omterm{prop:b2:biogeochemical-cycles:carbon}{الكسر المحمول جوًا} نحو الصفر إذ تلحق المصارف --- في
\omterm{thm:b2:biogeochemical-cycles:box}{نموذج الصندوق} الواحد؛ أما في الواقع فتتشبع المصارف ويهبط الكسر
أبطأ.
\textbf{7.} \qty{271}{GtC}.
\textbf{8.} $E^{*} = 10\times 100 = \qty{1000}{GtC}$: أي $285 +
1000/2.13 = \qty{755}{ppm}$.
\textbf{9.} $E(t) = 1000 + (271 - 1000)e^{-t/100}$؛ وعند $t = 80$ يكون
$E = 1000 - 729\times 0.449 = \qty{672}{GtC}$: أي \qty{600}{ppm}.
\textbf{10.} $Q = 10(1 - t/50)$: والحل الخاص $E_p = at + b$
عند $a = -20$ و$b = 3000$ (من $a = 10 - b/\tau$ و$0 = -0.2 -
a/\tau$)؛ فيكون $E = 3000 - 20t + (271 - 3000)e^{-t/100}$. وعند $t = 50$:
$2000 - 2729\times 0.607 = \qty{344}{GtC}$، أي \qty{447}{ppm}. ثم يخبو $E$
حرًّا: فعند $t = 80$ يكون $344\,e^{-0.3} = \qty{255}{GtC}$، أي
\qty{405}{ppm}.
\textbf{11.} \qty{600}{ppm} في مقابل \qty{405}{ppm}: فالسيناريو الثاني
يعيد الهواء، بحلول 2100، إلى ما هو عليه اليوم تقريبًا.
\textbf{12.} $\tau$ الواحد يجعل المصارف تمضي في العمل بالمعدل
نفسه على فائض يتقلص؛ أما في الواقع فالمصارف السريعة (الطبقة
الممتزجة، والغابات النامية) قد أخذت ما تستطيع
والإزالة الباقية بامتزاج المحيط العميق البطيء
والتجوية عند $\tau$ بالقرون إلى آلاف السنين، فيكون الهبوط
بعد الصفر الصافي أبطأ بكثير --- إذ يبقى التركيز ثابتًا
تقريبًا قرونًا بدل أن يهبط.
\textbf{13.} المحصول \qty{90}{kg}؛ والمنزوعة نترتته \qty{37.5}{kg}؛ والمغسول
\qty{22.5}{kg} من النتروجين في الهكتار في السنة.
\textbf{14.} $\qty{22.5}{kg}/\qty{14}{g/mol} = 1600$~مول نتروجين؛ والكربون
$1600\times 106/16 = \num{10600}$~مول، أي \qty{128}{kg}.
\textbf{15.} \num{10600}~مول من $\mathrm{O_2}$، أي \qty{340}{kg}.
\textbf{16.} $\num{10600}/0.25 = \num{42000}$~\unit{m^3}: فغسل
هكتار واحد يستطيع نزع الأكسجين من أربعين ألف متر
مكعب من الماء القاعي كل سنة --- والمسيسيبي يصرّف
مئة مليون هكتار.
\textbf{17.} $37.5\times 0.01 = \qty{0.375}{kg}$ من النتروجين بصيغة
$\mathrm{N_2O}$، أي $0.375\times 44/28 = \qty{0.59}{kg}$ من $\mathrm{N_2O}$؛
و$\times 270 = \qty{160}{kg}$ من مكافئ $\mathrm{CO_2}$.
\textbf{18.} $150\times 4 = \qty{600}{kg}$ من $\mathrm{CO_2}$: فالتصنيع
يفوق $\mathrm{N_2O}$ الحقل بأربعة إلى واحد، والاثنان
معًا ثلاثة أرباع الطن في الهكتار في السنة.
\textbf{19.} $\tau = 10$ سنوات؛ و$M^{*} = 5\times 10 = \qty{50}{t}$؛
والتركيز $50/(5\times 10^{7}) = \qty{1}{g/m^3} =
\qty{1000}{mg/m^3}$.
\textbf{20.} ثلاثون ضعف العتبة: أي إتخام شديد.
\textbf{21.} $1/\tau = 0.1 + 0.2 = 0.3$: أي $\tau = 3.3$ سنة؛ و$M^{*} =
5/0.3 = \qty{16.7}{t}$، أي \qty{330}{mg/m^3}.
\textbf{22.} $M^{*} = 1/0.3 = \qty{3.3}{t}$، أي \qty{67}{mg/m^3} --- وهي ما تزال
متخمة؛ ويُقترب منها إلى حدود \qty{10}{\%} بعد $\tau\ln 10 = 7.7$ سنة.
\textbf{23.} الرسوب ليس مصرفًا أحادي الاتجاه: فالفوسفور المخزون
فيه خلال سنوات الإتخام يُطلق راجعًا تحت الشروط
اللاأكسجينية التي يخلقها الإتخام نفسه (إذ ينحل الفوسفات
المرتبط بالحديد حين يُرجع الحديد)، فتغذي البحيرة
نفسها سنين بعد قطع الدخل الخارجي --- وهو
التحميل الداخلي، أي \omterm{def:b2:biogeochemical-cycles:reservoir}{مستودع} ثانٍ طويل \omterm{def:b2:biogeochemical-cycles:reservoir}{زمن المكث}.
\textbf{24.} $\qty{5000}{kg}/\qty{31}{g/mol} = 1.6\times 10^{5}$~مول فوسفور؛
والكربون $\times 106 = 1.7\times 10^{7}$~مول، أي \qty{205}{t} من الكربون في
السنة --- أي \qty{20}{g/m^2} على سطح البحيرة إذا كان عمقها
\qty{5}{m}، وهو ازدهار ثقيل.
\textbf{25.} \omterm{prop:b2:biogeochemical-cycles:carbon}{الكسر المحمول جوًا} $0.39$ على مدى 1850--2020؛ وسنة 2100،
\qty{600}{ppm} مع إبقاء الانطلاقات عند \qty{10}{GtC/yr} و
\qty{405}{ppm} مع قطعها إلى الصفر بحلول 2070 (وهو نموذج متفائل
ذو $\tau$ واحد)؛ والنتروجين المغسول \qty{22.5}{kg} في الهكتار
في السنة.
\end{solution}
